\chapter*{\sffamily Acknowledgements}
\addcontentsline{toc}{chapter}{Acknowledgements}

First and foremost, I wish to express my deepest gratitude to my supervisor, Andrés Marino Álvarez-Meza, and my co-supervisor, David Cárdenas-Peña. Their guidance, knowledge, patience, and unwavering support were fundamental throughout this journey. This research, as well as the personal and academic growth that accompanied it, would not have been possible without them.

I am profoundly grateful to Andrés for his guidance, his confidence in my abilities, and the academic rigor with which he accompanied every stage of this process. His advice and experience continually encouraged me to grow as a researcher, to face challenges with determination, and to always strive for excellence.

I extend a particularly heartfelt acknowledgement to David, who believed in me from the very beginning and remained by my side throughout this journey. Thank you for your constant presence, for listening to me, for supporting me during the most difficult moments, and for allowing me to experience this process with sensitivity and authenticity. Your trust, empathy, and encouragement gave me strength whenever the path became overwhelming. Beyond your invaluable academic guidance, your humanity left a profound and lasting mark on my life.

To my parents and my brother, I offer all my love and my deepest gratitude. Everything I have achieved rests upon the values, sacrifices, and unconditional love with which you have always supported me. Thank you for believing in me, for encouraging me to continue even when the path seemed uncertain, and for being the foundation that has sustained me throughout every stage of my life. Thank you for celebrating my achievements, accompanying me through my difficulties, and always reminding me of where I come from and everything I am capable of achieving. This accomplishment is also yours.

To my beloved fiancé, José, thank you for walking beside me throughout this journey. Thank you for your patience, your love, and your unwavering support; for comforting me on the difficult days, celebrating every small victory with me, and reminding me of my strength whenever I began to doubt it. Your companionship made the weight of this journey lighter and transformed its most meaningful moments into even more special memories. I am deeply grateful to have shared this chapter of my life with you and to know that we will continue building all the chapters that follow together.

I am especially grateful to all the members of the Control and Digital Signal Processing Group (GCPDS) at Universidad Nacional de Colombia, Manizales Campus, and the Automatics Research Group at Universidad Tecnológica de Pereira. The opportunities for learning, collaboration, and academic discussion that I found within these groups deeply enriched my academic development and allowed me to feel supported throughout this process.

I extend a particularly special acknowledgement to my best friend, María Isabel, who supported me constantly and remained by my side during some of the most difficult stages of this journey. Thank you for listening to me, for holding me up when I felt I could no longer continue, and for always giving me your affection, understanding, and unconditional friendship. Your companionship was a refuge during moments of uncertainty and a source of joy amid the difficulties. Having you by my side made this process far more bearable and meaningful.

I also express my sincere gratitude to Julián Pastrana, whose academic discussions, ideas, and perspectives contributed significantly to the development of this research and to my growth as a researcher. Thank you for sharing your knowledge with me, for the conversations that helped clarify so many questions, and for your companionship, which made my doctoral journey much more enjoyable and enriching.

Finally, I gratefully acknowledge the support provided by the project ``Alianza Científica con Enfoque Comunitario para Mitigar Brechas de Atención y Manejo de Trastornos Mentales Relacionados con Impulsividad en Colombia -- ACEMATE,'' through Grant No.~111091991908. Its support was essential for the development and completion of this research.
